\section{Introduction}
Cryptanalysis of block ciphers has long been a cornerstone of cryptographic research, driven by the need to ensure the security of data across various platforms. Block ciphers are essential in securing communication in environments ranging from the Internet of Things (IoT) and RFID systems to Smart Cards and sensor networks. The security requirements for these platforms differ significantly, often necessitating the use of lightweight cryptographic algorithms that can operate efficiently under constraints like limited area, energy, and code size. As a result, traditional ciphers such as the Advanced Encryption Standard (AES) may not be suitable for all applications, highlighting the importance of designing and thoroughly analyzing alternative cryptographic algorithms. 

Differential, linear and related key are among the most fundamental techniques in the cryptanalyst's toolkit. These methods have been playing important role in evaluating of the security of block ciphers. However, the complexity and time-consuming nature of these analyses pose significant challenges. 

The improvements in the area of cryptanalysis have led to the development of automated and semi-automated tools designed to enhance the efficiency of cryptanalysis. Among these tools, Mixed-Integer Linear Programming (MILP) has emerged as a particularly powerful method for cryptanalysis. MILP offers a flexible and effective approach to modeling and analyzing the security of block ciphers, making it easier to evaluate their resistance to cryptographic attacks.



\textbf{Organization of the paper}. The rest of the paper is organized as follows. Section 2 introduces Linear and mixed-integer linear programming with example. Section 3 defines ITUbee algorithm. Section 4 studies how to model block cipher components and. Section 6 introduces our MILP models for differantial, linear, relatekey differential cryptanalysis